\chapter{PENUTUP}

Pada bab terakhir ini dituliskan kesimpulan dan saran-saran yang dapat diambil berdasarkan materi-materi yang telah dibahas pada bab-bab sebelumnya.

\section{Kesimpulan}
Berdasarkan pembahasan yang telah dilakukan sebelumnya, penulis menarik beberapa kesimpulan berikut:
\begin{enumerate}
    \item Masalah \textit{electric bike-sharing system} dapat dimodelkan dengan masalah stokastik dua tahap.
    \item Permasalahan-permasalahan optimisasi seperti \textit{electric bike-sharing system} dapat diselesaikan dengan algoritma \textit{metaheuristic} yang salah satunya yaitu \textit{Fire Hawk Optimizer}. Simulasi \textit{electric bike-sharing system} di kota Yogyakarta dengan parameter-parameter yang telah dipaparkan sebelumnya memberikan gambaran bahwa \textit{electric bike-sharing system} ini dapat terealisasi dengan nilai keuntungan sebesar 274.030.876 dengan stasiun-stasiun yang dibangun yaitu, Sidomoyo, Monjali, Kentungan, Seturan, Tugu, Cik di Tiro, Malioboro, Simpang Empat Pelem, Ngastiharjo, Mergangsan, Rejowinangun, Bandara Adisucipto serta 55 sepeda listrik yang dibeli dan 80 trip yang diterima.
\end{enumerate}

\section{Saran}
Formulasi \textit{two-stage stochastic} pada permasalahan \textit{electric bike-sharing system} dan penyelesaiannya dengan menggunakan algoritma \textit{Fire Hawk Optimizer} serta beberapa kesimpulan yang telah dituliskan di atas,  penulis ingin menyampaikan beberapa saran sebagai berikut.
\begin{enumerate}
    \item Penelitian selanjutnya dapat dikembangkan dengan membuat formulasi yang lebih efisien sehingga proses komputasi tidak memakan daya ataupun waktu yang lama.
    \item Perlu dilakukan penelitian lebih lanjut terutama \textit{benchmarking} metode penyelesaian permasalahan \textit{electric bike-sharing system} seperti dengan algoritma \textit{metaheuristic} lain misal \textit{particle swarm optimization, simulated annealing}, algoritma genetika, ataupun yang lain.
\end{enumerate}